\documentclass[a4paper,12pt,twoside]{memoir}

% Castellano
\usepackage[spanish,es-tabla]{babel}
\selectlanguage{spanish}
\usepackage[utf8]{inputenc}
\usepackage[T1]{fontenc}
\usepackage{lmodern} % Scalable font
\usepackage{microtype}
\usepackage{placeins}

\RequirePackage{booktabs}
\RequirePackage[table]{xcolor}
\RequirePackage{xtab}
\RequirePackage{multirow}

% Links
\PassOptionsToPackage{hyphens}{url}\usepackage[colorlinks]{hyperref}
\hypersetup{
	allcolors = {red}
}

% Ecuaciones
\usepackage{amsmath}

% Rutas de fichero / paquete
\newcommand{\ruta}[1]{{\sffamily #1}}

% Párrafos
\nonzeroparskip

% Huérfanas y viudas
\widowpenalty100000
\clubpenalty100000

% Imágenes

% Comando para insertar una imagen en un lugar concreto.
% Los parámetros son:
% 1 --> Ruta absoluta/relativa de la figura
% 2 --> Texto a pie de figura
% 3 --> Tamaño en tanto por uno relativo al ancho de página
\usepackage{graphicx}
\newcommand{\imagen}[3]{
	\begin{figure}[!h]
		\centering
		\includegraphics[width=#3\textwidth]{#1}
		\caption{#2}\label{fig:#1}
	\end{figure}
	\FloatBarrier
}

% Comando para insertar una imagen sin posición.
% Los parámetros son:
% 1 --> Ruta absoluta/relativa de la figura
% 2 --> Texto a pie de figura
% 3 --> Tamaño en tanto por uno relativo al ancho de página
\newcommand{\imagenflotante}[3]{
	\begin{figure}
		\centering
		\includegraphics[width=#3\textwidth]{#1}
		\caption{#2}\label{fig:#1}
	\end{figure}
}

% El comando \figura nos permite insertar figuras comodamente, y utilizando
% siempre el mismo formato. Los parametros son:
% 1 --> Porcentaje del ancho de página que ocupará la figura (de 0 a 1)
% 2 --> Fichero de la imagen
% 3 --> Texto a pie de imagen
% 4 --> Etiqueta (label) para referencias
% 5 --> Opciones que queramos pasarle al \includegraphics
% 6 --> Opciones de posicionamiento a pasarle a \begin{figure}
\newcommand{\figuraConPosicion}[6]{%
  \setlength{\anchoFloat}{#1\textwidth}%
  \addtolength{\anchoFloat}{-4\fboxsep}%
  \setlength{\anchoFigura}{\anchoFloat}%
  \begin{figure}[#6]
    \begin{center}%
      \Ovalbox{%
        \begin{minipage}{\anchoFloat}%
          \begin{center}%
            \includegraphics[width=\anchoFigura,#5]{#2}%
            \caption{#3}%
            \label{#4}%
          \end{center}%
        \end{minipage}
      }%
    \end{center}%
  \end{figure}%
}

%
% Comando para incluir imágenes en formato apaisado (sin marco).
\newcommand{\figuraApaisadaSinMarco}[5]{%
  \begin{figure}%
    \begin{center}%
    \includegraphics[angle=90,height=#1\textheight,#5]{#2}%
    \caption{#3}%
    \label{#4}%
    \end{center}%
  \end{figure}%
}
% Para las tablas
\newcommand{\otoprule}{\midrule [\heavyrulewidth]}
%
% Nuevo comando para tablas pequeñas (menos de una página).
\newcommand{\tablaSmall}[5]{%
 \begin{table}
  \begin{center}
   \rowcolors {2}{gray!35}{}
   \begin{tabular}{#2}
    \toprule
    #4
    \otoprule
    #5
    \bottomrule
   \end{tabular}
   \caption{#1}
   \label{tabla:#3}
  \end{center}
 \end{table}
}

%
% Nuevo comando para tablas pequeñas (menos de una página).
\newcommand{\tablaSmallSinColores}[5]{%
 \begin{table}[H]
  \begin{center}
   \begin{tabular}{#2}
    \toprule
    #4
    \otoprule
    #5
    \bottomrule
   \end{tabular}
   \caption{#1}
   \label{tabla:#3}
  \end{center}
 \end{table}
}

\newcommand{\tablaApaisadaSmall}[5]{%
\begin{landscape}
  \begin{table}
   \begin{center}
    \rowcolors {2}{gray!35}{}
    \begin{tabular}{#2}
     \toprule
     #4
     \otoprule
     #5
     \bottomrule
    \end{tabular}
    \caption{#1}
    \label{tabla:#3}
   \end{center}
  \end{table}
\end{landscape}
}

%
% Nuevo comando para tablas grandes con cabecera y filas alternas coloreadas en gris.
\newcommand{\tabla}[6]{%
  \begin{center}
    \tablefirsthead{
      \toprule
      #5
      \otoprule
    }
    \tablehead{
      \multicolumn{#3}{l}{\small\sl continúa desde la página anterior}\\
      \toprule
      #5
      \otoprule
    }
    \tabletail{
      \hline
      \multicolumn{#3}{r}{\small\sl continúa en la página siguiente}\\
    }
    \tablelasttail{
      \hline
    }
    \bottomcaption{#1}
    \rowcolors {2}{gray!35}{}
    \begin{xtabular}{#2}
      #6
      \bottomrule
    \end{xtabular}
    \label{tabla:#4}
  \end{center}
}

%
% Nuevo comando para tablas grandes con cabecera.
\newcommand{\tablaSinColores}[6]{%
  \begin{center}
    \tablefirsthead{
      \toprule
      #5
      \otoprule
    }
    \tablehead{
      \multicolumn{#3}{l}{\small\sl continúa desde la página anterior}\\
      \toprule
      #5
      \otoprule
    }
    \tabletail{
      \hline
      \multicolumn{#3}{r}{\small\sl continúa en la página siguiente}\\
    }
    \tablelasttail{
      \hline
    }
    \bottomcaption{#1}
    \begin{xtabular}{#2}
      #6
      \bottomrule
    \end{xtabular}
    \label{tabla:#4}
  \end{center}
}

%
% Nuevo comando para tablas grandes sin cabecera.
\newcommand{\tablaSinCabecera}[5]{%
  \begin{center}
    \tablefirsthead{
      \toprule
    }
    \tablehead{
      \multicolumn{#3}{l}{\small\sl continúa desde la página anterior}\\
      \hline
    }
    \tabletail{
      \hline
      \multicolumn{#3}{r}{\small\sl continúa en la página siguiente}\\
    }
    \tablelasttail{
      \hline
    }
    \bottomcaption{#1}
  \begin{xtabular}{#2}
    #5
   \bottomrule
  \end{xtabular}
  \label{tabla:#4}
  \end{center}
}



\definecolor{cgoLight}{HTML}{EEEEEE}
\definecolor{cgoExtralight}{HTML}{FFFFFF}

%
% Nuevo comando para tablas grandes sin cabecera.
\newcommand{\tablaSinCabeceraConBandas}[5]{%
  \begin{center}
    \tablefirsthead{
      \toprule
    }
    \tablehead{
      \multicolumn{#3}{l}{\small\sl continúa desde la página anterior}\\
      \hline
    }
    \tabletail{
      \hline
      \multicolumn{#3}{r}{\small\sl continúa en la página siguiente}\\
    }
    \tablelasttail{
      \hline
    }
    \bottomcaption{#1}
    \rowcolors[]{1}{cgoExtralight}{cgoLight}

  \begin{xtabular}{#2}
    #5
   \bottomrule
  \end{xtabular}
  \label{tabla:#4}
  \end{center}
}



\graphicspath{ {./img/} }

% Capítulos
\chapterstyle{bianchi}
\newcommand{\capitulo}[2]{
	\setcounter{chapter}{#1}
	\setcounter{section}{0}
	\setcounter{figure}{0}
	\setcounter{table}{0}
	\chapter*{\thechapter.\enskip #2}
	\addcontentsline{toc}{chapter}{\thechapter.\enskip #2}
	\markboth{#2}{#2}
}

% Apéndices
\renewcommand{\appendixname}{Apéndice}
\renewcommand*\cftappendixname{\appendixname}

\newcommand{\apendice}[1]{
	%\renewcommand{\thechapter}{A}
	\chapter{#1}
}

\renewcommand*\cftappendixname{\appendixname\ }

% Formato de portada
\makeatletter
\usepackage{xcolor}
\newcommand{\tutor}[1]{\def\@tutor{#1}}
\newcommand{\course}[1]{\def\@course{#1}}
\definecolor{cpardoBox}{HTML}{E6E6FF}
\def\maketitle{
  \null
  \thispagestyle{empty}
  % Cabecera ----------------
\noindent\includegraphics[width=\textwidth]{cabecera}\vspace{1cm}%
  \vfill
  % Título proyecto y escudo informática ----------------
  \colorbox{cpardoBox}{%
    \begin{minipage}{.8\textwidth}
      \vspace{.5cm}\Large
      \begin{center}
      \textbf{TFG del Grado en Ingeniería Informática}\vspace{.6cm}\\
      \textbf{\LARGE\@title{}}
      \end{center}
      \vspace{.2cm}
    \end{minipage}

  }%
  \hfill\begin{minipage}{.20\textwidth}
    \includegraphics[width=\textwidth]{escudoInfor}
  \end{minipage}
  \vfill
  % Datos de alumno, curso y tutores ------------------
  \begin{center}%
  {%
    \noindent\LARGE
    Presentado por \@author{}\\ 
    en Universidad de Burgos --- \@date{}\\
    Tutor: \@tutor{}\\
  }%
  \end{center}%
  \null
  \cleardoublepage
  }
\makeatother

\newcommand{\nombre}{Rafael Martín Guerrero} %%% cambio de comando
\newcommand{\dni}{80087034E} %%% cambio de comando
\newcommand{\nomtutor}{Pedro Renedo Fernández} %%% cambio de comando
\newcommand{\nomcotutor}{} %%% cambio de comando
\newcommand{\titulo}{GCA} %%% cambio de comando
\newcommand{\departamentotutor}{Ingeniería Informática} %%% cambio de comando
\newcommand{\area}{Lenguajes y Sistemas Informáticos} %%% cambio de comando

% Datos de portada
\title{\titulo}
\author{\nombre}
\tutor{\nomtutor}
\date{\today}

\begin{document}

\maketitle


\newpage\null\thispagestyle{empty}\newpage


%%%%%%%%%%%%%%%%%%%%%%%%%%%%%%%%%%%%%%%%%%%%%%%%%%%%%%%%%%%%%%%%%%%%%%%%%%%%%%%%%%%%%%%%
\thispagestyle{empty}


\noindent\includegraphics[width=\textwidth]{cabecera}\vspace{1cm}

\noindent D. nombre tutor, profesor del departamento de nombre departamento, área de nombre área.

\noindent Expone:

\noindent Que el alumno D. \nombre, con DNI \dni, ha realizado el Trabajo final de Grado en Ingeniería Informática titulado \titulo.

\noindent Y que dicho trabajo ha sido realizado por el alumno bajo la dirección del que suscribe, en virtud de lo cual se autoriza su presentación y defensa.

\begin{center} %\large
En Burgos, {\large \today}
\end{center}

\vfill\vfill\vfill

% Author and supervisor
\begin{minipage}{0.45\textwidth}
\begin{flushleft} %\large
Vº. Bº. del Tutor:\\[2cm]
D. nombre tutor
\end{flushleft}
\end{minipage}
\hfill
\begin{minipage}{0.45\textwidth}
\begin{flushleft} %\large
Vº. Bº. del co-tutor:\\[2cm]
D. nombre co-tutor
\end{flushleft}
\end{minipage}
\hfill

\vfill

% para casos con solo un tutor comentar lo anterior
% y descomentar lo siguiente
%Vº. Bº. del Tutor:\\[2cm]
%D. nombre tutor


\newpage\null\thispagestyle{empty}\newpage




\frontmatter

% Abstract en castellano
\renewcommand*\abstractname{Resumen}
\begin{abstract}

Este proyecto tiene como objetivo desarrollar e implementar una plataforma web y un agente que permitan la gestión
centralizada de la configuración de red en los equipos de las aulas.La solución permite la segmentación de los
dispositivos en grupos y la aplicación de restricciones personalizadas para el acceso a recursos internos y externos,
incluyendo Internet.Se busca optimizar el control y administración de la red educativa, mejorando la seguridad y la
eficiencia en la gestión de los permisos de acceso.
\end{abstract}

\renewcommand*\abstractname{Descriptores}
\begin{abstract}
    Plataforma web, configuración de red, administración educativa, agente, control, segmentación, aplicación java
\end{abstract}

\clearpage

% Abstract en inglés
\renewcommand*\abstractname{Abstract}
\begin{abstract}
The objective of this project is to develop and implement a web platform and an agent that allows the centralized
management of the network configuration in the classrooms.The solution allows the segmentation of devices into groups
and the application of customized restrictions for access to internal and external resources, including the Internet.
The aim is to optimize the control and administration of the educational network, improving security and efficiency in
the management of access permissions.
\end{abstract}

\renewcommand*\abstractname{Keywords}
\begin{abstract}
    Web platform, network configuration, educational management, agent, control, segmentation, java application
\end{abstract}

\clearpage

% Indices
\tableofcontents

\clearpage

\listoffigures

\clearpage

\listoftables
\clearpage

\mainmatter
\capitulo{1}{Introducción}

La Universidad de Burgos (UBU) cuenta con más de 1000 equipos informáticos distribuidos en aulas entre 8 y 10 centros.
Cada una de estas aulas dispone de dispositivos destinados a satisfacer una serie de necesidades distintas, como por ejemplo,
la realización de operaciones matemáticas complejas, redacción de documentos, edición de contenido multimedia, entre otras.

Sin embargo, un mismo equipo no está destinado a un único uso y con las aplicaciones instaladas pueden permitir el acceso a otros
recursos y dispositivos no deseados, esto supone un riesgo para la seguridad de la información y puede derivar en un uso
indebido, como acceder a información no autorizada durante una prueba de evaluación.

En los últimos años, se ha observado un incremento de las peticiones entre los docentes, para limitar los recursos a
los que se pueden acceder desde los ordenadores.
Por este motivo, surge la necesidad de implementar una gestión centralizada en las aulas, que permita, de una manera sencilla, controlar
los accesos de los distintos equipos, atendiendo a características diferentes, como nombres de equipo, aulas, centros\ldots

De acuerdo con Microsoft\cite{microsoftQuAccess}existen cuatro formas de control de acceso:
    \begin{itemize}
        \item Control de accesso dicrecional (DAC): un propietario concede acceso a otros usuarios.
        \item Control de acceso basado en roles (RBAC): se otorga acceso según funciones empresariales definidas.
        \item Control de acceso obligatorio (MAC): se concede acceso en forma de autorización regulado por una autoridad central.
        \item Control de acceso basado en atributos (ABAC): se da de manera flexible en función de variables y condiciones.
    \end{itemize}

Actualmente, los usuarios de los equipos, mediante el protocolo 802.1X\cite{eswiki:802.1X} han de ingresar sus credenciales
y tras una autentificación satisfactoria, se les asigna una dirección IP, que les permite hacer uso del dispositivo, navegar
por Internet, acceder a los distintos recursos de red\ldots.

Existen herramientas que pueden ayudar a limitar el acceso a los recursos, como software específico, no obstante,
estad opciones pueden resultar costosas.
También se pueden configurar uno a uno los permisos de cada dispositivo, pero ello requiere
que manualmente, un administrador tenga que cambiar los distintos permisos, y llevar un control de cada uno de ellos.

La plataforma propuesta, busca optimizar costes y ofrecer una solución práctica, diseñada a medida para la UBU.
A diferencia de otras propuestas, esta iniciativa se enfoca en resolver problemas reales sin añadir funciones innecesarias
que acaban dificultando el uso.
Para su implementación, solo requiere un servidor web encargado de gestionar las configuraciones y una puesta a punto inicial
en cada equipo para su registro en el sistema.
Posteriormente, un software se encargará de aplicar las configuraciones
necesarias en cada dispositivo de forma automática.

\capitulo{2}{Objetivos del proyecto}

Este apartado explica de forma precisa y concisa cuales son los objetivos que se persiguen con la realización del proyecto. Se puede distinguir entre los objetivos marcados por los requisitos del software a construir y los objetivos de carácter técnico que plantea a la hora de llevar a la práctica el proyecto.

\capitulo{3}{Conceptos teóricos}

En aquellos proyectos que necesiten para su comprensión y desarrollo de unos conceptos teóricos de una determinada materia o de un determinado dominio de conocimiento, debe existir un apartado que sintetice dichos conceptos.

Algunos conceptos teóricos de \LaTeX{} \footnote{Créditos a los proyectos de Álvaro López Cantero: Configurador de Presupuestos y Roberto Izquierdo Amo: PLQuiz}.

\section{Secciones}

Las secciones se incluyen con el comando section.

\subsection{Subsecciones}

Además de secciones tenemos subsecciones.

\subsubsection{Subsubsecciones}

Y subsecciones. 


\section{Referencias}

Las referencias se incluyen en el texto usando cite~\cite{wiki:latex}. Para citar webs, artículos o libros~\cite{koza92}, si se desean citar más de uno en el mismo lugar~\cite{bortolot2005, koza92}.


\section{Imágenes}

Se pueden incluir imágenes con los comandos standard de \LaTeX, pero esta plantilla dispone de comandos propios como por ejemplo el siguiente:

\imagen{escudoInfor}{Autómata para una expresión vacía}{.5}



\section{Listas de items}

Existen tres posibilidades:

\begin{itemize}
	\item primer item.
	\item segundo item.
\end{itemize}

\begin{enumerate}
	\item primer item.
	\item segundo item.
\end{enumerate}

\begin{description}
	\item[Primer item] más información sobre el primer item.
	\item[Segundo item] más información sobre el segundo item.
\end{description}
	
\begin{itemize}
\item 
\end{itemize}

\section{Tablas}

Igualmente se pueden usar los comandos específicos de \LaTeX o bien usar alguno de los comandos de la plantilla.

\tablaSmall{Herramientas y tecnologías utilizadas en cada parte del proyecto}{l c c c c}{herramientasportipodeuso}
{ \multicolumn{1}{l}{Herramientas} & App AngularJS & API REST & BD & Memoria \\}{ 
HTML5 & X & & &\\
CSS3 & X & & &\\
BOOTSTRAP & X & & &\\
JavaScript & X & & &\\
AngularJS & X & & &\\
Bower & X & & &\\
PHP & & X & &\\
Karma + Jasmine & X & & &\\
Slim framework & & X & &\\
Idiorm & & X & &\\
Composer & & X & &\\
JSON & X & X & &\\
PhpStorm & X & X & &\\
MySQL & & & X &\\
PhpMyAdmin & & & X &\\
Git + BitBucket & X & X & X & X\\
Mik\TeX{} & & & & X\\
\TeX{}Maker & & & & X\\
Astah & & & & X\\
Balsamiq Mockups & X & & &\\
VersionOne & X & X & X & X\\
} 

\capitulo{4}{Técnicas y herramientas}

\section{Metodologías}\label{sec:metodologias}

\subsection{Scrum}\label{subsec:scrum}

La metodología ágil utilizada ha sido Scrum\cite{scrumorg:whatis}, mediante la cual se ha organizado el trabajo en
interaciones cortas llamadas sprints.
En cada uno de estos se establecen una serie de tareas, ordenadas por prioridad.

Este enfoque permite una mayor flexibilidad, ya que no se necesita tener todo definido desde el inicio del proyecto,
si surgen imprevistos, como cambios de requisitos, nuevas ideas, o cualquier tipo de problema, se introducen nuevas tareas
en los siguientes sprints.

Este marco de trabajo encaja muy bien en equipos pequeños que trabajan de manera autónoma, favoreciendo una autogestión
eficiente y una mejora continua del proyecto.

\subsection{Test Driven Development (TDD)}\label{subsec:tdd}

En la misma línea de metodologías ágiles, se ha optado por seguir la práctica de TDD\cite{browserstack:tdd}, en la cual
se escriben los tests antes de implementar el código que se necesita probar.

TDD consta de tres fases repetitivas:

\begin{enumerate}
    \item \textbf{Fase roja :} se escriben los test destinados a comprobar una determinada funcionalidad.
    Inicialmente, estos tests fallan.
    \item \textbf{Fase verde :} se implementa el mínimo código necesario para que los tests anteriores pasen.
    \item \textbf{Refactorización :} se optimiza el código a la vez que se asegura que todos los test siguen siendo correctos.
\end{enumerate}

\section{Gestión del proyecto}\label{sec:gestion}

\subsection{Control de versión}\label{subsec:control}

Para la gestión de versiones del proyecto se ha escogido Git\cite{git:official}, una herramienta de control de versiones distribuido.

Permite llevar un seguimiento de todos los cambios realizados en el código, además de facilitar la creación de ramas
independientes que pueden usarse para implementar nuevas funcionalidades o corregir errores.

Existen otras herramientas de gestión como Subversion (SVN)\cite{svn:official}, la cual se había considerado también:

\begin{itemize}
    \item \textbf{Modelo:} Git es un sistema distribuido, mientras que Subversion (SVN) es centralizado.
    \item \textbf{Repositorio local:} Git almacena todo el historial del proyecto localmente; SVN solo guarda una copia de trabajo.
    \item \textbf{Velocidad:} Git realiza operaciones locales mucho más rápido que SVN, que depende del servidor.
    \item \textbf{Gestión de ramas:} Git facilita la creación y fusión de ramas de forma ligera; en SVN, este proceso es más pesado y complejo.
    \item \textbf{Integración con plataformas:} Git se integra de forma nativa con plataformas como GitHub o GitLab, mientras que SVN tiene una integración más limitada.
    \item \textbf{Soporte y comunidad:} Git cuenta con una comunidad más amplia, con mayor soporte y documentación actualizada que SVN.
\end{itemize}

\subsection{Repositorio}\label{subsec:repositorio}

Las opciones contempladas fueron GitHub o Gitlab, debido a su integración nativa con Git.
Finalmente se ha optado por utilizar GitHub\cite{wikipedia:github}, ya que se trata de una plataforma muy popular en proyectos open source, es
muy intuitiva y está enfocada a desarrolladores individuales y equipos colaborativos.

Por su parte, Gitlab está más enfocada a entornos de empresas, es más completo pero más complejo de utilizar.

Además, GitHub incluye una serie de herramientas de integración CI / CD llamda GitHub Actions que se utilizará en
conjunto con otras herramientas que describiremos en otras secciones.

\subsection{Plataforma de gestión}\label{subsec:plataforma}

Existen infinidad de plataformas enfocadas a la gestión de proyectos; Trello, Jira, Asana, ZenHub \ldots

Hemos optado por utilizar ZenHub\cite{zenhub:gettingstarted}, ya que es una herramienta que se integra fácilmente con GitHub, permite crear
sprints, épicas y tableros Kanban.

Si bien es cierto que la herramienta más utilizada por excelencia es Jira, no obstante considero que tiene demasiadas
funcionalidades que no van a ser aprovechadas en este proyecto, por lo que es mas conveniente utilizar algo menos potente.

\subsection{Comunicación}\label{subsec:comunicacion}

Para cualquier tipo de comunicación, se utilizará email y Teams.

\section{Lengujes de programación}\label{sec:lenguajes}

\subsection{Cliente}\label{subsec:cliente}

Para el desarrollo de la aplicación cliente, se debía escoger un lenguaje de scripts o cercano a estos, ya que el objetivo
era construir una herramienta capaz de obtener configuraciones y aplicarlas a los dispositivos de forma sencilla.

Entre los candidatos estaban PowerShell, Java, y Python.
Todos son lenguajes que permiten satisfacer los requisitos del proyecto, con ligeras diferencias:
\begin{itemize}
    \item \textbf{PowerShell:} aunque es potente, limita su uso a dispositivos con sistema operativo Windows,
    restringiendo así la portabilidad.
    \item \textbf{Java:} gracias a su arquitectura multiplataforma, permite ejecutar programas en cualquier entorno.
    No obstante, generar un archivo ejecutable (.exe) que funcione universalmente resulta complejo, y en caso de
    distribuir un archivo .jar, se requeriría que los dispositivos clientes tuvieran Java instalado previamente.
    \item \textbf{Python:} en cambio, ofrece una solución más práctica: mediante sus librerías, es posible generar
    ejecutables de forma sencilla, compatibles con distintos sistemas operativos, y sin necesidad de instalar
    software adicional en los dispositivos cliente.
\end{itemize}

Por esto, se decide utilizar Python\cite{python:docs} para desarrollar la aplicación de escritorio.

\subsection{API}\label{subsec:api}

En cuanto al desarrollo de APIs, los dos cantidatos por excelencia son Spring Boot y Node.js.

Node.js ha adquirido gran relevancia en los ultimos años, ya que permite crear APIs de una manera muy sencilla,
siendo especialmente útil en aplicaciones que precisan de una gran velocidad y eficiencia en tiempo real.
Por otro lado, Spring Boot, ofrece un amplio conjunto de módulos, que permiten configurar la aplicación atendiendo
a las necesidades de cada proyecto, a cambio de una mayor complejidad de desarrollo.

Se ha escogido Spring Boot\cite{spring:docs} debido a que, aunque el proyecto inicial es relativamente pequeño, ofrece la  posibilidad
de escalar la aplicación y gestionar un mayor número de usuarios en un futuro.
Además, integra módulos (beans, core, security\ldots) que facilitan la implementación de mecanismos avanzados de desarrollo,
como el control de accesos de usuarios.

\subsection{Frontal}\label{subsec:frontal}

Finalmente, para el frontend, se evaluaron dos alternativas: Angular o React.
React proporciona gran flexibilidad y es ideal para proyectos en los que se quiere tener control total de la arquitectura
y librerias usadas, aunque depende en gran medida de librerias de terceros para su desarrollo.
Angular, en cambio ofrece un framework completo, con una gran cantidad de funcionalidades predefinidas, lo que facilita
el mantenimiento y la optimización del código, especialmente en proyectos grandes.

Se ha optado por Angular\cite{angular:docs}, priorizando su robustez, escalabilidad y la gran cantidad de herramientas que proporciona.

\section{Arquitectura}\label{sec:arquitectura}

\subsection{Patrón de diseños}\label{subsec:patron}

En este proyecto se han aplicado distintos patrones de diseño, adaptados a las diferentes necesidades.

La aplicación de escritorio del cliente, desarrollada en Python, emplea una estructura basada en la gestión dinámica
de una única ventana, ocultando y mostrando distintos elementos.
Esto es similar al concepto de SPA (Single Page Application) en el ámbito web.

La API REST desarrollada en Spring Boot, sigue una arquitectura MVC (Modelo-Vista-Controlador)\cite{wikipedia:mvc}.

El frontend en Angular adopta un patrón similar a MVC, donde los componentes representan la vista, y el controlador
son los servicios que encapsulan el modelo de datos y la comunicación con la API.

\section{Herramientas de desarrollo}\label{sec:desarrollo}

\subsection{Cliente}\label{subsec:desarrollo_cliente}

Para el desarrollo de la aplicación cliente se ha utilizado el Entorno de Desarrollo Integrado (IDE) PyCharm\cite{pycharm:official},
ya que se encuentra plenamente enfocado a Python, ofrece multitud de herramientas y permite realizar configuraciones rápidas.

Otra opción era Visual Studio Code, pero para este lenguaje en concreto, se necesitan descargar multitud de extensiones
para poder desarrollar cómodamente.

Para la instalación de librerias externas, se utiliza PIP\cite{pip:official}, un gestor de paquetes que permite
instalar, actualizar y gestionar librerías.

\subsection{API}\label{subsec:desarrollo_api}

Para nuestra API con Spring Boot, se ha optado por el IDE IntelliJ IDEA\cite{intellij:official} frente a Eclipse,
por la sencilla razón de que IntelliJ es mucho más potente, intuitivo y dispone de todas las herramientas necesarias,
sin necesidad de instalar nada.

En conjunto con este IDE, se ha añadido una herramienta de automatización de compilación moderna y potente, como es
Gradle\cite{gradle:official}, permite compilar, construir, testear y desplegar aplicaciones de software.

En combinación con Spring Boot, destaca su uso frente al de Maven, ya que es más ligero, flexible y permite tener un control de
los despliegues.

\subsection{Frontend}\label{subsec:desarrollo_frontend}

El frontend se ha desarrollado utilizando Visual Studio Code\cite{vscode:official}, también es posible utilizar IntelliJ y tener ambos
proyectos (API y Frontend) en el mismo IDE, sin embargo, no está tan optimizado para la parte front y la sensación de
programar el frontal dista mucho de ser óptima.

Además, utilizamos NPM (Node Package Manager)\cite{npm:official} como gestor de paquetes para instalar librerias
externas de una manera simple y segura.

\section{Integración contínua}\label{sec:integracion}

\subsection{Integración y despliegue continuos (IC/DC)}\label{subsec:integracion_continua}

Para automatizar compilaciones, pruebas y despliegue, se ha utilizado GitHub actions\cite{github:actions}, en este caso, al
escoger GitHub como repo, viene incluido sin necesidad de instalar nada más.

\subsection{Cobertura y calidad de código}\label{subsec:cobertura}

SonarQube Cloud\cite{sonarcloud:official} se ha seleccionado como plataforma para analizar la calidad y cobertura del código fuente.
La razón de su uso se debe a que no es necesario instalar ningún servidor, basta con conectarlo al respositorio de GitHub y
se puede consultar todos los resultados online.

\subsection{Despliegue}\label{subsec:despliegue}

Docker\cite{docker:official} es una plataforma que permite empaquetar una aplicación y todas sus dependencias en un
contenedor portable y que puede ejecutarse en cualquier entorno

\section{Documentación}\label{sec:documentacion}

La documentación se ha desarrollado utilizando LaTeX\cite{latex:official}, en lugar de otras herramientas de procesamiento de texto como Word.
Word funciona muy bien para documentos simples y además es fácil e intuitivo desde el principio.
No obstante, LaTeX permite un control de todo el formato, genera documentos de una calidad tipográfica
excelente y es ideal para documentos técnicos.
Sin embargo, es necesario un periodo de adaptación para aprender a usarlo.

Como editor LaTeX se utiliza IntelliJ Idea junto con el plugin TeXiFy-IDEA y el compilador pdfLaTeX.

Visualizamos los documentos PDF generados con SumatraPDF por integrarse de manera excelente con las herramientas
mencionadas.


\section{Librerias}\label{sec:librerias}

En este apartado mencionamos las librerias más importantes, sin contar aquellas que se encuentran preinstaladas en los
frameworks utilizados.

\subsection{Cliente}\label{subsec:librerias_cliente}

\textbf{PyInstaller}

PyInstaller\cite{pyinstaller:official} es una herramienta que permite convertir programas escritos en Python en ejecutables independientes para
Windows, Linux y macOS.


\textbf{Inno Setup Compiler}

Inno Setup\cite{innosetup:official} es una herramienta gratuita para crear instaladores para aplicaciones de Windows.
Permite empaquetar archivos, configuraciones y recursos necesarios, generando un asistente de instalación (.exe) que
facilita la distribución y el despliegue de programas.

\subsection{API}\label{subsec:librerias_api}


\textbf{Jjwt-api}

Jjwt-api forma parte de la biblioteca JJWT (Java JWT)\cite{jjwt:official}, una de las librerías más populares para trabajar con
JSON Web Tokens (JWT) en aplicaciones Java.


\textbf{JUnit}

JUnit\cite{junit:official} es uno de los frameworks de testing más populares para aplicaciones Java.
Permite crear y ejecutar de manera estructurada tests unitarios, ayudando a garantizar que cada componente de un
programa funcione correctamente de forma aislada.


\subsection{Frontend}\label{subsec:librerias_frontend}



\textbf{Angular CLI}

Angular CLI (Command Line Interface)\cite{angularcli:official} es una herramienta oficial para crear, desarrollar,
construir y mantener aplicaciones Angular desde la línea de comandos.


\textbf{Tailwind CSS}

Tailwind CSS\cite{tailwind:official} es un framework de CSS utilitario que permite diseñar interfaces de usuario
rápidamente mediante la combinación de pequeñas clases directamente en el HTML .


\textbf{DaisyUI}

DaisyUI\cite{daisyui:official} es un complemento de componentes UI construido sobre Tailwind CSS .

Ofrece componentes predefinidos como botones, modales, menús, tabs, alertas, etc., todos diseñados usando clases de Tailwind.




\capitulo{5}{Aspectos relevantes del desarrollo del proyecto}

En este punto se detallan los aspectos más relevantes en el desarrollo del proyecto, desde la primera reunión con el tutor para la
elección del tema, hasta la entrega del mismo.

\section{Elección del proyecto}\label{sec:eleccion-del-proyecto}

En la primera reunión con el tutor, se enumeraron los posibles temas para el Trabajo de Fin de Grado.

El motivo de la elección de este tema viene motivado por experiencias de desarrollo pasadas, donde realicé una aplicación
web para la gestión de las escuelas deportivas del centro donde cursé el CFGS de Desarrollo de Aplicaciones Web (DAW).

\section{Métodos de trabajo}\label{sec:metodos-de-trabajo}

A pesar de no contar con un equipo de trabajo, se decidió seguir una metodología ágil, concretamente Scrum; al no
contar, evidentemente, con un equipo de trabajo, no se ha aplicado de una manera estricta, pero se han seguido los
principios básicos de la misma:
\begin{itemize}
    \item \textbf{Sprint}: se han realizado sprints de dos semanas, donde se planifica el trabajo a realizar
    según un orden de prioridades.
    \item \textbf{Reuniones semanales}: cada semana, los lunes concretamente, se mantuvieron reuniones con el tutor
    para revisar el avance del proyecto (a modo de \textit{demo}) y resolver dudas.
    \item \textbf{Backlog}: en esta lista se han ido añadiendo, durante el curso del desarrollo del proyecto,
    todo tipo de tareas a realizar (de análisis, documentación,
    diseño, desarrollo, pruebas\ldots).
\end{itemize}

En cuanto al desarrollo con Scrum, no se define un conjunto de fases específicas, sino que se trabaja de una forma
iterativa e incremental: en cada sprint se van añadiendo nuevas funcionalidades y se van corrigiendo errores.

\section{Fase de análisis}\label{sec:fase-de-analisis}

Durante este periodo, se dedicaron varias sesiones a la recopilación de información sobre el tema a desarrollar, las distintas
necesidades de los usuarios y la manera de satisfacerlas.

La idea principal era crear una herramienta que permitiese gestionar, de forma centralizada, el acceso a la red de
los dispositivos que dispone la Universidad.
Para ello, se estudió la forma de gestionar este acceso desde la consola de comandos y se descubrieron herramientas que
permitían esta gestión, como el firewall de Windows y Netsh.

Uno de los principales problemas que surgieron fue la manera de identificar a cada dispositivo de forma única.
La primera idea fue utilizar el GUID del dispositivo, pero se descubrió que los dispositivos del centro son
clonados, y, por lo tanto, comparten el mismo, por lo que se tuvo que descartar irremediablemente esta vía.

Ante esta dificultad, se planteó la posibilidad de combinar este identificador GUID con el nombre del dispositivo para
crear una huella única.

Estos dispositivos se guardarían en una base de datos junto con las configuraciones que se pretendan aplicar.
Un detalle importante que se especificó en esta fase, es que contaríamos con grupos o localizaciones de dispositivos
que formarían una jerarquía y, así, se podrían aplicar configuraciones de carácter general a todos los dispositivos que
pertenecen a cada grupo.

\section{Fase de diseño}\label{sec:fase-de-diseno}

Barajamos la posibilidad de utilizar un script que se mantuviese en ejecución en segundo plano, pero se
descartó por ser poco eficiente y seguro.
En su lugar, decidimos utilizar un ejecutable y mantenerlo activo de forma continua, mediante el programador de
tareas de Windows; de esta manera, se restringe la posibilidad de que un usuario pueda modificar el script y ejecutar
comandos que puedan comprometer la seguridad del sistema.
Las tecnologías propuestas para el ejecutable fueron Java y Python, pero la creación de un ejecutable (.exe) en Java
causó muchos problemas y no se pudo realizar.
Con un archivo JAR sí se pudo crear, pero generaba la necesidad de tener instalado Java en cada dispositivo.
Por eso se optó por utilizar Python, ya que permite crear un ejecutable independiente del sistema operativo.

En cuanto a la base de datos (BD) se decidió utilizar un modelo relacional, porque los datos a almacenar son de
carácter estructurado y van a tener relación entre ellos.

Para la API se decidió utilizar un modelo REST, ya que es el más utilizado en la actualidad y permite una fácil
interacción entre el cliente y el servidor.
La herramienta elegida fue Spring Boot y Java, debido a que está enfocada a la creación de servicios REST y cuenta con una
gran cantidad de herramientas y librerías ya instaladas.

Por último, se decidió crear una interfaz web para realizar toda la gestión de manera centralizada: por compatibilidad
y por ser uno de los frameworks actuales más utilizados, se decidió utilizar Angular.

Con toda esta estructura, aseguramos que la aplicación sea escalable y fácil de mantener, e incluso si se quisiese
crear una aplicación para smartphone en un futuro, se podría reutilizar la parte backend, agilizando el proceso de
desarrollo de la app.

\section{Fase de desarrollo}\label{sec:fase-de-desarrollo}

La construcción del código comienza con el ejecutable en Python: una primera versión con una ventana
y un formulario para iniciar sesión, y, tras la autenticación, un botón para registrar el dispositivo
en la base de datos.

El principal desafío consistió en crear un instalador.
Para ello, se utilizó Inno Setup combinado con el ejecutable (.exe) generado mediante la librería PyInstaller; esto
permitió no solo instalar la aplicación en el dispositivo, sino también mantenerla en ejecución en segundo plano,
mostrando un icono en la barra de tareas.


El siguiente paso fue la creación de la API con la base de datos, que se realizó de forma paralela al ejecutable.
La configuración inicial de la API con Spring Boot fue lo más tedioso, porque fue necesario instalar todas aquellas
dependencias que se van a utilizar, además de establecer la estructura de carpetas y archivos.
La forma de autenticación elegida fue mediante JWT, una de las más utilizadas en la actualidad, lo que hace posible
una fácil interacción entre el cliente y el servidor.
Se implementaron, de manera gradual, los distintos casos de uso de la API, generando diferentes endpoints para
cada uno de estos.
Añadimos SonarCloud para el análisis de la calidad del código con GitHub Actions; sin embargo, inicialmente, se
desactivaron los tests porque la BD no es en memoria ni se contaba con Docker, además tuvimos que utilizar una versión
de Java inferior a la 21, ya que SonarCloud soporta hasta la 17.

Para el desarrollo de la interfaz web, fue necesario descargar Node.js y Angular CLI con tal de generar el proyecto
y sus dependencias.
A partir de este paso, cualquier librería externa se añade mediante el gestor de paquetes NPM\@.
Utilizamos Tailwind CSS + DaisyUI para el diseño de la interfaz, ya que incluye estilos predefinidos para
maquetar la aplicación.

La interfaz web se desarrolla de forma paralela a la API: primero se escriben los tests de la API, y a continuación, se escribe
el código necesario para que pasen, aplicando el desarrollo dirigido por pruebas (TDD); por último, se implementa la interfaz web.
Una de las funcionalidades más difíciles de implementar fue la estructura arbórea de los grupos situada en el menú
izquierdo de la interfaz; tenía que cargar de forma dinámica el listado de dispositivos que pertenecen a cada uno en
la parte derecha, además de mantener la concordancia con los datos.
Esto es clave, porque facilita al usuario la gestión de los dispositivos de una manera intuitiva y rápida.

La gestión de plantillas y comandos se implementa de una manera en la que se crean los comandos ``fijos``, y estos
son añadidos a las distintas plantillas.
No obstante, se descubrió que esta lógica no era muy eficaz.
Si queríamos, por ejemplo, añadir un valor a un comando de una dirección IP, teníamos que crear un comando para cada
dirección, lo que resultaría tremendamente tedioso y se perdería la agilidad que pretendemos conseguir de forma general con GCA.
Para solventar este problema, se estableció una estructura básica para añadir valores tal que el valor de un comando que contenga \verb|{{variable}}|
sirve para cambiar el valor de esa variable al asignarse a una plantilla.

Cuando terminamos con la construcción de la API (o, al menos, el desarrollo de aquellas estructuras que se consideraron
necesarias para su funcionamiento), se decidió mejorar la aplicación Python para el cliente.
Inicialmente, había demasiada lógica en cada archivo .py, por lo que creamos clases más pequeñas con una lógica bien definida.
El funcionamiento de mostrar / ocultar botones limita el desarrollo y añade lógica innecesaria, por lo que se opta por
crear frames o páginas y añadir en estas los elementos necesarios.
En nuestro caso, creamos una página para iniciar sesión y otra para mostrar el panel de usuario una vez iniciada.
Otro punto a destacar aquí es que, cuando intentamos realizar el registro de un equipo y aplicar la configuración,
no se sabía qué estaba ocurriendo, más allá de alguna alerta programada.
Introdujimos un recuadro para ir mostrando todos los logs de la aplicación mediante diferentes colores: los de
información en azul, advertencias en naranja y errores en rojo.

Tras conseguir una primera aplicación de la configuración asignada a un dispositivo, comenzaron los problemas:
\begin{itemize}
    \item Al deshabilitar la red para el dispositivo, se perdía la conectividad a la API.
    \item Tras aplicar la configuración de plantilla, quitar la asignación y añadir otra diferente, no se aplicaba
    correctamente.
    \item Al bloquear toda conexión saliente de la interfaz, no podíamos conectar de nuevo a la API.
\end{itemize}

Tuvimos que trabajar en el algoritmo para aplicar la configuración; así, el funcionamiento queda definido de la siguiente
manera:

\begin{enumerate}
    \item Se restaura la copia de seguridad.
    \item Se eliminan los registros de plantillas anteriores.
    \item A partir de aquí, se comprueba de manera periódica la existencia de configuración a aplicar.
    \item Se obtiene el valor de la huella digital almacenada en el dispositivo.
    \item Si existe una huella digital, se solicita al servidor la configuración a aplicar.
    \item Si se reciben datos de configuración:
    \begin{itemize}
        \item Se comprueba si la configuración obtenida es más actual que la que tiene el equipo.
        \item Si hay cambios, se restaura la copia de seguridad y se aplican los comandos.
        \item Se actualizan los valores en el registro del dispositivo.
    \end{itemize}
    \item Si no hay configuración o está vacía, se revisa si se ha eliminado la asignación del equipo y restaura.
    \item Se espera el tiempo definido antes de reiniciar el ciclo.
\end{enumerate}

La copia de seguridad se realiza una única vez en el momento de registrar el equipo, de forma que nos aseguramos de que
el dispositivo pueda volver a su estado original ante cualquier imprevisto.
Esta copia incluye la configuración del firewall y la configuración de red.
En el proceso de restauración también se habilitan de nuevo las interfaces de red y se configuran las reglas de firewall
para mantener la conexión a las URL's configuradas en el json como permitidas.

\section{Fase de pruebas}\label{sec:fase-de-pruebas}

Las pruebas se realizan de manera continua a lo largo de todo el ciclo del desarrollo: para el backend utilizamos las
librerías JUnit para pruebas unitarias, y Spring Boot Test para pruebas de integración.

En primer lugar, escogimos la misma BD con la que inicia la API para ejecutar los tests, pero esto generaba
conflictos con los datos, por lo que se decidió utilizar una BD en memoria, H2, para realizar las pruebas.
Para la configuración de la BD en memoria, creamos un nuevo archivo de application.properties dentro del directorio
de test/resources.
Este archivo contiene la configuración de la BD en memoria, y se carga automáticamente al ejecutar los tests.
Aquí surgió un problema: los scripts se ejecutaban antes de que Hibernate crease la BD dando lugar a errores de inserción.
Para que funcionase correctamente, según lo propuesto en StackOverflow\cite{stackoverflowSpringBootError},
fue necesario añadir la siguiente propiedad spring.jpa.defer-datasource-initialization=true para invertir el orden, y
que primero se generase la estructura y después se insertase los datos.

Una vez configurada la BD en memoria, se vuelven a activar los tests en SonarQube Cloud, pero se producen una serie
de errores relacionados con el contexto de Spring.
Aquellos test que hacen uso de la anotación @SpringBootTest levantan el contexto, permiten inyectar beans y utilizar
todos los componentes, y como resultado, se ejecutaban de nuevo los scripts de la base de datos.
Existía una confusión entre las tablas generadas por JPA y los scripts que se ejecutaban para la BD en memoria,
por lo que se ajustó poniendo todos los nombres de tablas en minúsculas, y para los tests de spring,
estableciendo el perfil activo a la hora de ejecutar los tests de ''tests''.
De esta manera, solamente se ejecutan aquellas propiedades dentro de la carpeta test/.

Mediante este enfoque conseguimos un entorno de pruebas aislado, reproducible y compatible con SonarQube.

Las pruebas de la aplicación cliente se hicieron de manera manual, tratando de cubrir todos aquellos casos que
se puedan dar, como por ejemplo:
\begin{itemize}
    \item Asignar una plantilla de configuración vacía.
    \item Asignar una plantilla a un dispositivo y al grupo directo al que pertenece y ver qué plantilla se obtiene.
    \item Eliminar la plantilla asignada y ver qué sucede con la configuración.
    \item Reiniciar / Apagar el dispositivo y comprobar que se aplique la configuración correcta.
\end{itemize}

De esta manera, aseguramos que la funcionalidad más importante esté cubierta.

Por último, para la aplicación web, al alojar un número reducido de funcionalidades, simplemente se comprueba que
toda la gestión de dispositivos, grupos, comandos y plantillas sea correcta, es decir, que se pueda añadir, eliminar,
modificar y listar.
\capitulo{6}{Trabajos relacionados}

La gestión centralizada de dispositivos, es una práctica habitual en la administración de sistemas,
y casi desde sus inicios, se han desarrollado herramientas para la gestión de dispositivos de red.
Esta gestión ha evolucionado, y sigue evolucionando hacia una automatización de tareas, la orquestación de servicios y
la gestión de configuraciones.

En este apartado se exponen aquellos artículos científicos y herramientas existentes que guardan similitud con
el presente trabajo.

\section{Articulos cientificos}\label{sec:articulos-cientificos}

\textbf{An Application to Automate the Configuration Management and Assessment of Practically Based Learning Outcomes in Computer Networking }

Este estudio describe una aplicación desarrollada para automatizar la configuración y evaluación de escenarios prácticos
en laboratorios de redes\cite{malliarakis2019application}.
La herramienta permite aplicar configuraciones a dispositivos reales y simulados, facilitando  la evaluación de
resultados de aprendizaje prácticos.
Aunque se centra en el ámbito educativo, su enfoque en la automatización y gestión de configuraciones es aplicable a
entornos uniersitarios más amplios.

\textbf{Configuration Management and Security }

Este artículo\cite{bellovin2001configuration} analiza los desafíos de la gestión de configuración en sistemas de gran
escala, destacando la importancia de políticas centralizadas para garantizar la seguridad y el cumplimiento normativo.
Aunque no se enfoca específicamente en la educación, se describe la necesidad de herramientas que permitan la gestión
de configuraciones de manera eficiente y segura.

\section{Herramientas}\label{sec:herramientas}

\textbf{System Center Configuration Manager (SCCM)}

SCCM\cite{microsoftConfigMgr} es una solución de Micrososft para la gestión centralizada de configuraciones y
actualizaciones.
Permite aplicar configuraciones a grupos de dispositivos, gestionar actualizaciones de software \ldots
Aunque es una herramienta comercial, su enfoque en la gestión de configuraciones es similar al
propuesto en este trabajo.
Para su uso, es necesario tener instalado además Microsoft AD para permitir gestionar dispositivos y usuarios dentro
de la red.

\textbf{Active Directory (AD) or Entra ID)}

AD\cite{microsoftADDS} y Entra ID\cite{microsoftEntra} (anteriormente conocido como Azure AD) son soluciones de Microsoft
para la gestión centralizada mediante servicios que controlan el acceso a recursos y la configuración de dispositivos
en entornos empresariales.
La diferencia principal entre ambos es que Entra ID está diseñado para entornos en la nube, mientras que AD se utiliza
principalmente en entornos locales.

\textbf{Netsh}

Netsh\cite{microsoftNetsh} es una herramienta de comandos de Microsoft que permite la configuración y gestión
de dispositivos de red en sistemas operativos Windows, ya sea de forma local o remota.

 \begin{table}[H]
        \centering
        \begin{tabular}{lcccc}
            \toprule
            \textbf{Características} & \textbf{GCA} & \textbf{SCCM} & \textbf{AD} & \textbf{Netsh} \\
            \midrule
            Automatización de configuraciones & \cellcolor{verde}Sí & \cellcolor{verde}Sí & \cellcolor{verde}Sí & \cellcolor{rojo}No \\
            Gestión centralizada              & \cellcolor{verde}Sí & \cellcolor{verde}Sí & \cellcolor{verde}Sí & \cellcolor{rojo}No \\
            Priorización por dispositivos     & \cellcolor{verde}Si & \cellcolor{amarillo}Limitada & \cellcolor{verde}Sí & \cellcolor{rojo}No \\
            Interfaz de usuario               & \cellcolor{verde}Sí & \cellcolor{verde}Sí & \cellcolor{verde}Sí & \cellcolor{rojo}No (terminal) \\
            Curva de aprendizaje              & \cellcolor{verde}Baja & \cellcolor{rojo}Alta & \cellcolor{rojo}Alta & \cellcolor{amarillo}Media \\
            Dependencia con otras herramientas & \cellcolor{verde}No & \cellcolor{rojo}Sí & \cellcolor{rojo}Sí & \cellcolor{verde}No \\
            Multiplataforma                   & \cellcolor{verde}Sí* & \cellcolor{amarillo}Limitada & \cellcolor{amarillo}Limitada & \cellcolor{rojo}No \\
            \bottomrule
        \end{tabular}
        \caption{Comparativa de las herramientas de los proyectos con codificación de colores.}
        \label{tab:comparativa-herramientas-proyectos}
 \end{table}

Esta comparativa revela que, aunque soluciones como SCCM y Active Directory ofrecen capacidades avanzadas de gestión de
configuraciones y escalabilidad, requieren de hardware específico, son económicamente costosas y están diseñadas para
entornos empresariales.
En contraste, la propuesta de GCA se centra en la simplicidad y accesibilidad, permitiendo su uso en entornos educativos
y de investigación sin la necesidad de infraestructura costosa y compleja, pudiendo ser implementada en cualquier
tipo entorno, si así se desea.

No obstante, la elección de la herramienta adecuada dependerá de las necesidades específicas de cada entorno y de los
recursos disponibles.

\section{Fortalezas y debilidades del proyecto}\label{sec:fortalezas-y-debilidades-del-proyecto}

\begin{table}[ht]
    \centering
    \begin{tabular}{@{}p{7cm} p{7cm}@{}}
        \toprule
        \textbf{Fortalezas} & \textbf{Debilidades} \\
        \midrule
        Automatización de tareas repetitivas & Dependencia de la conectividad \\
        Gestión centralizada y escalable & Complejidad en la implementación \\
        Priorización de configuraciones & Riesgo de obsolescencia tecnológica \\
        Integración con herramientas existentes & Requiere primera instalación manual \\
        Mejora en la seguridad & Posibles problemas de compatibilidad \\
        Compatible con cualquier sistema operativo & Necesidad de mantenimiento continuo \\
        \bottomrule
    \end{tabular}
    \caption{Fortalezas y Debilidades del Proyecto}
    \label{tab:fortalezas-debilidades}
\end{table}
\capitulo{7}{Conclusiones y Líneas de trabajo futuras}

Todo proyecto debe incluir las conclusiones que se derivan de su desarrollo. Éstas pueden ser de diferente índole, dependiendo de la tipología del proyecto, pero normalmente van a estar presentes un conjunto de conclusiones relacionadas con los resultados del proyecto y un conjunto de conclusiones técnicas. 
Además, resulta muy útil realizar un informe crítico indicando cómo se puede mejorar el proyecto, o cómo se puede continuar trabajando en la línea del proyecto realizado. 



\bibliographystyle{plain}
\bibliography{bibliografia}

\end{document}
